\documentclass{situationist-preprint}

\title{The Madagascar Case}
\msid{LNG-02-b}
\msfamily{LNG-02}
\msdiscipline{Philosophy of Language}
\mstype{Companion manuscript}
\mskeywords{reference; proper names; causal theory; Madagascar; Evans; reference change}

\begin{document}
\maketitle

\begin{abstract}
Evans's Madagascar case is the standard counterexample to the pure causal theory of names: \enquote{Madagascar} once named a part of the African mainland and now names the island, although the practice of its use descends from the earlier one. LNG-02-a uses the case as one of three fixed points for its hybrid theory. We examine the case at length. We locate the break in the chain at Marco Polo's link, explain why the causal component fails there, show that the descriptive component determines the island, and derive from the Transmission lemma that all later uses refer to the island. We then answer the objection that multiple grounding, without any descriptive component, accounts for the case.
\end{abstract}

\section{The case}

According to the story Evans reports, \enquote{Madagascar} was a name used by local speakers for a portion of the African mainland; Marco Polo, misunderstanding his informants, took it to be the name of the great island, and his usage spread until the name came to refer to the island \citep{Evans1973}. Whatever the historical details, the case as described has three features. The present practice descends by acquisition from the earlier one. The earlier practice was grounded in the mainland. The present practice refers to the island. A pure causal theory, on which a name refers to whatever lies at the origin of the chain of its uses \citep{Kripke1980}, predicts that the name still refers to the mainland, and so gives the wrong verdict.

\section{Where the chain breaks}

On the account of LNG-02-a, a link is reference-preserving only if the acquiring speaker's subsequent uses are governed by the intention to refer to what the source referred to and are not sustained predominantly by information whose dominant causal source is some other object (Definition~3.2 there). Marco Polo acquired the name with the first condition satisfied: he intended to use it as his informants did. The second condition fails. His uses of the name in his account of his travels were sustained by information about an island: its size, its position off the coast, the reports of its inhabitants. The dominant causal source of that information was the island, not the mainland. So his link is not reference-preserving, and the practice that descends from him is not intact.

The causal component of the practice that includes Marco Polo's chain therefore fails to determine an object (Definition~3.3(a) of LNG-02-a). The descriptive component of that practice, the body of information its users predominantly associated with the name, is satisfied by exactly one object, the island (Definition~3.4). By Lemma~4.3 of LNG-02-a, any rule satisfying the constraints of that manuscript follows the descriptive component in this situation; the hybrid rule (Definition~3.5) assigns the island. Once the island is the referent at the stage of Marco Polo's readers, Corollary~4.7 applies: every later link being reference-preserving, every later use refers to the island. The present referent is thus explained by one descriptive takeover followed by ordinary causal transmission.

Two features of this diagnosis matter. First, the break is not a matter of the acquirer's intentions alone; Marco Polo intended to defer. What breaks the link is the source of the information on which his use depended. Second, the descriptive component does not fix reference for the earlier practice, in which the causal component succeeded: the local speakers referred to the mainland whatever they believed about it, as the G\"odel case requires (Lemma~4.1 of LNG-02-a).

\begin{objection}[Multiple grounding]
Devitt explains reference change without descriptions: after Marco Polo, the name was grounded again and again in perceptual encounters with the island, and these groundings came to outweigh the original ones \citep{Devitt1981}. The descriptive component is idle.
\end{objection}

\begin{reply}
Multiple grounding explains how the island reference became entrenched, not how it began. The later groundings were episodes in which speakers applied \enquote{Madagascar} to the island in perceptual contact with it, and they did so because they already took the name to refer to the island. That belief was transmitted from Marco Polo's account, which no one grounded perceptually in the island under the name before the account existed. If the first uses after the break referred to nothing, or to the mainland, the later groundings would be groundings of a new name that happened to share a form with the old. The hybrid account says that the first uses referred to the island by their descriptive component, so that the later groundings continue one practice. The two accounts are compatible; the descriptive component is needed at exactly the point where the causal component fails.
\end{reply}

\begin{objection}
On this account, a single traveler's confusion changes the reference of a name for a whole community. That makes reference too fragile.
\end{objection}

\begin{reply}
It does not change the reference of the local practice, which remains intact and continues to refer to the mainland. It starts a second practice. Corollary~4.7 of LNG-02-a guarantees stability once the new practice is under way, and stability within a practice is all that the robustness of reference requires.
\end{reply}

\section{Conclusion}

The Madagascar case is a break in the causal chain at a link whose uses were sustained by information from another object. The descriptive component then determines the island, and causal transmission carries that reference to the present. The case confirms that names refer causally, with a descriptive component that takes over when the chain breaks.

\begin{thebibliography}{9}\small
\bibitem[Devitt(1981)]{Devitt1981} Devitt, M. (1981). \emph{Designation}. Columbia University Press.
\bibitem[Evans(1973)]{Evans1973} Evans, G. (1973). The causal theory of names. \emph{Proceedings of the Aristotelian Society, Supplementary Volume}, 47, 187--208.
\bibitem[Kripke(1980)]{Kripke1980} Kripke, S.~A. (1980). \emph{Naming and Necessity}. Harvard University Press.
\end{thebibliography}

\end{document}
